\documentclass{article}

% Language setting
% Replace `english' with e.g. `spanish' to change the document language
\usepackage[spanish]{babel}
\usepackage[most]{tcolorbox}

% Set page size and margins
% Replace `letterpaper' with `a4paper' for UK/EU standard size
\usepackage[letterpaper,top=2cm,bottom=2cm,left=3cm,right=3cm,marginparwidth=1.75cm]{geometry}

% Useful packages
\usepackage{amsmath}
\usepackage{graphicx}
%\usepackage{subfigure}
\usepackage{caption}
\usepackage{subcaption}
\usepackage[colorlinks=true, allcolors=blue]{hyperref}

\newtcolorbox[auto counter, number within=section]{exercise}[1][]{
    colback=blue!5!white,
    colframe=blue!75!black,
    fonttitle=\bfseries,
    title=Ejercicio \thetcbcounter,
    #1 % Allows you to add custom options locally (like labels)
}




\title{Práctica 3. Reinforcement learning}
\author{Arturo Gil y Juan José Cabrera}

\begin{document}
\maketitle

\section{Introducción}
La mayor parte de los conceptos vistos hasta el momento en la asignatura se pueden clasificar como algoritmos de aprendizaje supervisado. Dentro de este tipo de aprendizaje consideramos algoritmos a los que se le proporciona un conjunto de datos de entrenamiento etiquetados. El rol del algoritmo es encontrar una función que permita hallar una mapeo lo más aproximado posible entre las entradas y las salidas. En el ámbito del aprendizaje por refuerzo (\emph{reinforcement learning}) se plantea un problema diferente y complejo, pues se plantea que un agente alcance algún grado de aprendizaje, aún cuando no existen datos de entrenamiento. Además, los datos de entrenamiento se generan mientras el agente intenta interactuar con el entorno.

A continuación indicamos algunos ejemplos de reinforcement learning:
\begin{itemize}
    \item Un robot en un laberinto (figura \ref{fig:robot_in_a_maze}): asuma que tenemos un robot en un laberinto. Se trata de un robot omnidireccional. Así pues, en cada momento el robot puede moverse en 4 direcciones (izquierda, derecha, arriba o abajo). Consideramos que en todo momento podemos conocer la posición del robot en el laberinto. El objetivo del problema es que el robot sea capaz de encontrar el camino óptimo que le lleve permita encontrar la salida, cuando el robot parte desde diferentes puntos de inicio. Así pues, el aprendizaje consiste en que el robot se forme una especie de mapa mental que le permita tomar la mejor acción (moverse a la izquierda, derecha, arriba abajo) en cada punto del mapa.
    \item Imaginemos ahora que tenemos un robot bípedo como el mostrado en la figura \ref{fig:robot_bipedo}. En este caso, nos gustaría poder programar al robot para que pudiera andar. Imaginemos también que tenemos un modelo dinámico que nos permite simular al sistema (figura \ref{fig:robot_bipedo_gymnasium}). Finalmente, sabemos que la marcha bípeda es extremadamente compleja, pues todas las variables del sistema están acopladas (cuando se le aplica un par a una articulación también se generan reacciones de fuerzas y pares en otras articulaciones). ¿Por qué no ponemos el simulador a probar diferentes soluciones para poder aprender un algoritmo de marcha? Durante la simulación cada vez que el robot realice un movimiento que le permita avanzar, esto se premiará. Por contra, cada vez que el robot se caiga al suelo, se castigará este hecho.
\end{itemize}

Merece la pena decir aquí que, como en casi todos los ámbitos del conocimiento, existen soluciones alternativas. Así pues, un ingeniero especialista en el ámbito de la robótica móvil nos diría que también se podría plantear la construcción de un mapa de ocupación con alguna técnica de \emph{mapping} o \emph{SLAM} para después utilizar el algoritmo $A^*$ y planificar una ruta óptima. En el caso de los robots andadores, también hay otras formas de plantear el problema. Típicamente, un ingeniero definiría con cuidado el \emph{gait} de cada articulación y después intentaría optimizarlo en base a los parámetros dinámicos del robot.

\textbf{Resumen de la práctica}: en esta práctica se recordarán algunos conceptos de aprendizaje por refuerzo y, en concreto, del algoritmo conocido como Q-learning. Seguidamente, se verán algunos ejemplos de Q-learning usando la librería Gymnasium para python. \textbf{El estudiante experimentará de diferentes formas con estos entornos y observará de qué forma afectan al aprendizaje.}

\begin{figure}
    \centering
    \includegraphics[width=0.7\linewidth]{figures/maze.png}
    \caption{Un robot en un laberinto}
    \label{fig:robot_in_a_maze}
\end{figure}

\begin{figure}
    \centering
    \begin{subfigure}{0.25\textwidth}
        \includegraphics[width=\linewidth]{figures/bipedal_tron.png}
        \caption{Un robot bípedo.}
        \label{fig:robot_bipedo}
    \end{subfigure}
    \hfill
    \begin{subfigure}{0.6\textwidth}
        \includegraphics[width=\linewidth]{figures/bipedal_gymnasium.png}
         \caption{Un robot bípedo simulado en Gymnasium.}
        \label{fig:robot_bipedo_gymnasium}
    \end{subfigure}
    \caption{Un robot bípedo real (arriba) y simulado (abajo).}
    \label{fig:robots_bipedos}
\end{figure}


\section{El repositorio de prácticas y Gymnasium}
Se os proporciona un repositorio de Github con algunos scripts que usan el algoritmo Q-learning y están basados en la librería Gymnasium. Gymnasium (\url{https://gymnasium.farama.org/}) es una librería para python que incluye una serie de entornos y agentes. Es especialmente adecuada para probar y desarrollar algoritmos de aprendizaje por refuerzo. Si quieres observar estos entornos, visita \url{https://gymnasium.farama.org/environments/classic_control/}, por ejemplo para ver los ejemplos de reinforcement learning para problemas clásicos de control (el péndulo invertido, cart-pole, etc...).


\section{Algunos conceptos de \textit{Reinforcement Learning}}

\subsection{Q-learning}
El algoritmo Q-learning es uno de los más famosos en el ámbito del aprendizaje por refuerzo. Se considera un algoritmo que no precisa de modelos del entorno ni del agente. El algoritmo busca la estimación de una tabla Q donde cada fila representa un estado $s$ y las columnas representan las posibles acciones $a$ que puede realizar el agente en cada estado. En cada celda de la tabla se almacena un valor $Q(s, a)$ que estima la recompensa futura (sobre todos los estados futuros hasta el final del episodio) de tomar la acción $a$ cuando el agente se encuentra en el estado $s$. La tabla se calcula durante el entrenamiento usando la ecuación de Bellman. Durante la fase de funcionamiento (test), el agente ejecuta la acción $a$ para la cual el valor en la tabla es máximo.
Asumimos que en un momento determinado de tiempo $t$, el estado del agente y del entorno se puede describir mediante $s_t \in S$ (es decir, un estado $s_t$ dentro del conjunto de posibles estados que describen al agente y al entorno). Asumimos que en ese mismo instante el agente decide realizar una acción $a_t \in A$. Contamos con un modelo de transición de estados (también llamado un modelo de movimiento). Este modelo se representa de dos formas diferentes:
\begin{itemize}
    \item En el caso más sencillo: $a_{t+1}=\delta(s_t, a_t)$. Dado que el robot se encuentra en $s_t$, si se ejecuta la acción $a_t$ entonces el estado cambiará a ser $s_{t+1}$. En el ejemplo inicial del laberinto, el robot se moverá a la siguiente posición, dependiendo de la acción $a_t$ que se realice (arriba, abajo, izquierda o derecha) y, además, dependiendo de si hay una pared en el laberinto o si el espacio está libre.
    \item En el caso más general, hablaremos de un modelo de transición de estado. Estos modelos se consideran ejemplos de MDP (\emph{Markov Decision Processes}) y se representan mediante una función de probabilidad
    $P(s_{t+1}|s_t, a_t)$. Es decir, dado que el agente y entorno se encuentran en $s_t$ y se ejecuta la acción $a_t$, podemos establecer una función de probabilidad sobre el estado futuro $s_{t+1}$.
\end{itemize} 

%Qué es el aprendizaje por refuerzo?
%¿Qué es el Q-learning?
%¿Qué es Gymnasium?
%Por qué



\subsection{Etapas y episodios}
Una etapa (paso o step) del algoritmo comprende el caso en el que el agente se encuentra en un estado $s_t$, realiza una acción $a_t$ que la conduce a otro estado $s_{t+1}$ (y recibe una recompensa $r_t$).
En algunos procesos podemos establecer de forma natural cuándo se inicia y se finaliza la misión. En el caso del robot en el laberinto, el episodio comienza con el robot en la casilla inicial y finaliza cuando el robot llega a la salida (o bien cuando se ha superado un número máximo de pasos). Realizada la última acción por el robot, la misión finaliza y, por tanto, podemos decir que el episodio acaba.

En otras situaciones, como en el caso en el que se quiere controlar un proceso continuo, el concepto de episodio no es adecuado, pues el control debe funcionar \textit{ad eternum}.

Normalmente, un episodio durará una serie de pasos máximos (por ejemplo, 200 pasos).
Ejecutaremos, también, un número de episodios con estados iniciales diferentes. 

\subsection{Entrenamiento}

El objetivo del entrenamiento es maximizar la suma de las recompensas de cada episodio. Esto supone que las recompensas deben estar bien definidas. Según se definan las recompensas, así será el resultado final.  Así pues, los pesos de las recompensas deben ser elegidos cuidadosamente y con rigor para que se produzca un correcto aprendizaje. En los casos de aprendizaje en los que se desea que un robot camine, se premia con +1 cada vez que el robot avanza en la dirección deseada y con -10, por ejemplo, cuando el robot pierde el equilibrio y cae.



\subsection{Train y test}
En el paradigma del aprendizaje, se plantean generalmente dos fases separadas: 
\begin{itemize}
    \item Se plantea inicialmente una etapa en la que el agente se entrena, obteniendo conocimiento sobre el entorno, según realiza acciones y obtiene recompensas.
    
\item Se plantea, después, otra etapa en la que se prueba el algoritmo .
\end{itemize}

 En ambas etapas nos interesa medir el comportamiento del agente para evaluar sus conocimientos. También es frecuente que tengamos que depurar el algoritmo para saber por qué no está aprendiendo (o bien, por qué está aprendiendo). Aquí se apunta alguna información que debe ser apuntada en cada episodio para observar el aprendizaje:
 \begin{itemize}
     \item Suma de recompensas en cada episodio.
     \item Histograma de los estados visitados en todos los episodios.
     \item Imagen de la tabla Q.
 \end{itemize}

En el segundo de los casos, lo más interesante sería poder comprobar si el agente ha visitado todos los estados posibles y, en cada uno de ellos, ha podido probar todas las acciones disponibles.



\subsection{Exploitation, exploration and reward}
Definimos a continuación los conceptos de explotación, exploración y recompensa (\textit{exploitation, exploration and reward}):
\begin{itemize}
    \item \textit{Exploitation}: hace referencia al cálculo de una acción a realizar cuando se hace uso del conocimiento actual sobre la acción que produce una mejor recompensa futura.
    \item \textit{Exploration}: se habla de exploración cuando el agente toma decisiones aleatorias. Esto permite llegar a estados que permiten mejorar el conocimiento que tenemos sobre las recompensas. Responde a la pregunta ¿qué pasa si estando en este estado llevo a cabo una acción aleatoria $a_t$, obtendré una recompensa futura mayor?
\end{itemize}

\textbf{En general, explotar el conocimiento que tenemos permite maximizar la recompensa actual, pero la exploración permite probar nuevas formas de resolver el problema que podrían producir una recompensa total futura mayor.}


%Ahora mismo es mejor que aprendáis a resolver este problema poco a poco, os llevará más tiempo y esfuerzo ahora, pero ahorraréis tiempo en el futuro. Si le preguntáis a la IA ahora, os dará una solución rápida pero, seguramente, no aprenderéis, y a la larga os costará más aplicar esta solución a un problema más complejo.

En un único paso del algoritmo no es posible explorar y explotar al mismo tiempo, pues únicamente podemos aplicar la acción una vez y, a continuación, nos encontraremos en otro estado. Así pues, en ocasiones se habla del conflicto entre la exploración y la explotación del conocimiento que se tiene en un determinado momento.

En relación con esto, en un momento determinado del ciclo de aprendizaje, nos podríamos preguntar, ¿cómo de seguros estamos de que tenemos un conocimiento correcto sobre la tabla $Q(s,a)$. Por ejemplo, imaginemos que el agente ha llegado a un estado $s_t=8$ y consulta la tabla $Q(s,a)$ donde se obtiene una acción $a_t$=GO RIGHT. Después de realizar esa acción, el agente obtiene una recompensa $r_t=1$. Imaginemos que ahora estamos en un episodio diferente y hemos llegado al mismo estado $s_t=8$ de nuevo. Nos preguntamos ahora si a) debemos repetir la acción anterior y asegurar una recompensa inmediata de 1 punto o, b) elegir otra acción diferente, esperando que esto dé lugar, a la larga, a una recompensa final mayor.

La forma más académica de resolver el dilema entre la exploración y la explotación es plantear que incertidumbre tenemos sobre la matriz Q. Imaginemos que el agente que está tomando decisiones es un estadístico. Podría evaluar las acciones que se habían tomado con anterioridad y evaluar la incertidumbre con la que se conoce esta información. Cuando la incertidumbre sobre las estimaciones sea  baja, el estadístico debería decidir que ya no es necesario más acciones exploratorias.

\subsection{Algoritmo greedy y epsilon greedy}
En esta práctica se plantea una forma sencilla de articular la exploración y la explotación. Esta forma de proceder se denomina $\epsilon$-greedy. Funciona de la siguiente manera, se inicializa $\epsilon=1$ y se programa una función decreciente en el tiempo de este parámetro:
\begin{verbatim}
    for episode in episodes: 
        p=normrnd()
        if p < epsilon:
           exploration
        else
           exploitation
        epsilon = t*epsilon
\end{verbatim}

Así pues, con probabilidad $\epsilon$ se explora y con probabilidad $1-\epsilon$ se explota el conocimiento. El valor de $\epsilon$ se va reduciendo en cada episodio, pues es preferible tomar acciones aleatorias durante los primeros episodios del aprendizaje y refinar el conocimiento durante los episodios finales. Esta estrategia se denomina $\epsilon$-greedy con decaimiento. La estrategia consiste en reducir $\epsilon$ según avanza la simulación al multiplicarlo por un factor $\delta$:
\begin{equation}
    \epsilon_{t+1}= \epsilon_{t}*\delta
\end{equation}
Esta operación se hace, bien durante cada paso de la simulación o bien en cada episodio. Además, se fija un valor mínimo:
\begin{equation}
    \epsilon_{t+1}= \max(\epsilon_{min}, \epsilon_{t}*\delta)
\end{equation}

Usaremos los siguientes valores:
\begin{itemize}
    \item $\delta=0.99$.
    \item $\epsilon_{0}=1.0$ (valor inicial).
    \item $\epsilon_{min}=0.01$
\end{itemize}

%Id est: durante los primeros episodios se explora más (en busca de nuevo conocimiento) y durante los últimos se explota más el conocimiento que se posee. Esta estrategia se denomina método epsilon-greedy de exploración-explotación.






\section{Un primer entorno}
Comenzaremos viendo un primer entorno sencillo y planteando una técnica para que el agente aprenda. El entorno se observa en la figura \ref{fig:cliff_walking_environment}. Se trata del entorno denominado \textit{Cliff Walking}. En este entorno, el agente (un gnomo):
\begin{itemize}
    \item Comienza siempre a moverse desde la posición [3, 0] de una rejilla de 4x12 y debe llegar al objetivo situado en la posición [3, 11].
    \item Si el agente llega a la posición deseada, el episodio finaliza.
    \item Hay un acantilado en todas las celdas [3, 1...10]. 
    \item Si el gnomo se mueve a una celda perteneciente al acantilado, volverá a la casilla de salida.
\end{itemize}

El objetivo es que el agente sea capaz de llegar desde la casilla de salida hasta el destino en el menor número de pasos de simulación posible. 

Este entorno es, en extremo sencillo. Fácilmente, podríamos escribir a mano la tabla $Q$ óptima que resolvería el problema, haciendo que el gnomo siguiera la secuencia obvia: Arriba, derecha, derecha, derecha, ..., abajo. No obstante, se desea comenzar con este entorno sencillo para capacitar al estudiante para resolver problemas más compplejos. De esta manera, se busca que el estudiantado:
\begin{itemize}
    \item Se familiarice con el algoritmo Q-learning cuando se intuye o conoce el resultado.
    \item Deduzca la ecuación de actualización de la tabla Q (la ecuación de Bellman) por prueba y error.    
\end{itemize}

Puedes consultar toda la documentación sobre este entorno aquí:  \url{https://gymnasium.farama.org/environments/toy_text/cliff_walking/}

\begin{figure}
    \centering
    \includegraphics[width=0.8\linewidth]{figures/cliff_walking.png}
    \caption{El entorno cliff-walking}
    \label{fig:cliff_walking_environment}
\end{figure}


\subsection{Estados, acciones y recompensas}
El entorno de Gymnasium codifica 3x12+1 estados diferentes, pues el agente puede estar en cualquiera de las celdas de las 3 filas superiores y, además, puede encontrarse en la casilla de salida. El agente no puede estar en ninguna de las celdas del acantilado ni en la celda de meta, pues esto ocasiona el final del episodio.

\subsubsection{Codificación del estado}
El estado del problema debe codificar la posición del gnomo en cada instante. Una solución obvia consistiría en definir el estado con las coordenadas $(i, j)$ de cada celda, recordando que en python ambas se numeran comenzado en 0. Sin embargo, resulta más conveniente codificar los estados con un único entero según se muestra en la figura \ref{fig:states}. Así pues, cuando el agente se encuentra en la posición $(i, j)$ del entorno, el estado entero se codifica como $s=i\times n_c+ j$, donde $n_c$ es el número de columnas del entorno (12). Por ejemplo, la posición de salida del agente se codifica con el entero 36 ($36 = 3\times 12+0$).

\begin{figure}
    \centering
    \includegraphics[width=0.8\linewidth]{figures/states.png}
    \caption{La codificación de estados del problema.}
    \label{fig:states}
\end{figure}


\subsubsection{Acciones}
En cada posición el agente puede realizar las siguientes acciones, codificada cada una de ellas con un número entero:
\begin{itemize}
    \item 0: Arriba.
    \item 1: Derecha.
    \item 2: Abajo.
   \item  3: Izquierda
\end{itemize}

\subsubsection{Recompensas}
El gnomo recibe una recompensa de -1 en cada paso. Se obtiene una recompensa de -100 cuando se cae por el acantilado. Se desea que el gnomo alcance la meta con una suma de recompensas en el episodio máxima.
\begin{exercise}
    ¿Cuál es la recompensa de un episodio en el que el gnomo completa la tarea de forma óptima?
\end{exercise}

\subsubsection{La tabla Q}
Con todo lo anterior, la tabla Q tiene la forma de esta matriz:
\begin{equation*}
    Q = \begin{bmatrix}
            q_{0,0} & q_{0,1} & q_{0,2} & q_{0,3}\\
            q_{1,0} & q_{1,1} & q_{1,2} & q_{1,3}\\
            \ldots & \ldots & \ldots & \ldots\\
            q_{35,0} & q_{35,1} & q_{35,2} & q_{35,3}\\
        \end{bmatrix}
\end{equation*}
La matriz tiene tantas filas como estados posibles se han codificado y tantas columnas como posibles acciones a realizar en cada estado. Cada uno de los elementos $q_{i,j}$ almacena el valor de $Q(i, j)$ que estima las recompensas futuras que se obtendrían al realizar la acción $j$ cuando el agente se encuentra en el estado $i$. La técnica Q-learning propone una manera de estimar esta matriz para maximizar la recompensa de las acciones realizadas por el agente. En resumen: debemos calcular cada elemento de la matriz \texttt{q\_table} para que \texttt{q\_table[s][a]} estime la recompensa futura (sobre todos los estados
futuros hasta el final del episodio).

\subsubsection{Un primer script}
Si abrimos el script \texttt{cliff\_q\_learning\_test.py}, encontraremos el siguiente script. Léalo con atención. La clase \texttt{QLearning} se proporciona programada parcialmente, para que el estudiante la complete. Importante:
\begin{itemize}
    \item En estos momentos la matriz Q no se conoce, con lo que el script genera una matriz Q aleatoria.
    \item Si se ejecuta el script, el agente no es capaz de resolver la tarea: en ocasiones no se moverá y en otras ocasiones repetirá los movimientos sin llegar a su objetivo.
\end{itemize}
\begin{verbatim}
from qlearning.qlearning import QLearning
import gymnasium as gym

def test_qlearning():
    environment = gym.make('CliffWalking-v1', render_mode="human")
    qlearning = QLearning(environment=environment)
    qlearning.load_random_q_table()
    #qlearning.load_q_table(filename='qtable.npy')
    print('Testing mode!')
    qlearning.test(total_episodes_test=1)

if __name__ == "__main__":
    test_qlearning()
\end{verbatim}

% Define a counter within sections


\begin{exercise}
Examine el método test de la clase QLearning, describa su funcionamiento.
\end{exercise}

\begin{exercise}
Ejecute el método test 50 veces. Describa las acciones que realiza el agente. ¿Se consigue completar el episodio?
\end{exercise}

\begin{exercise}
Calcule la recompensa media de cada episodio y plotee el resultado.
\end{exercise}

El Teorema del mono infinito afirma que un mono pulsando teclas al azar sobre un teclado durante un periodo de tiempo indefinido podrá escribir finalmente ``Hamlet'' de William Shakespeare (\url{https://es.wikipedia.org/wiki/Teorema_del_mono_infinito}). De la misma manera, podríamos generar tablas $Q$ aleatorias y pensar que alguna de ellas permitiría que el gnomo fuera desde la casilla de salida hasta la casilla de llegada. Si, por algún casual, ocurre esto, este caso sería fácil de detectar y nos permitiría aproximar la tabla Q (pues cada valor de Q(i,j) tendría el valor total de recompensa del episodio). Sin embargo, la tabla Q calculada así dará difícilmente una solución óptima del problema: con mucha suerte el gnomo llegará al final... y con muchísima suerte la trayectoria realizada será óptima. También podríamos pensar qué hacer con diferentes tablas Q que resuelvan el problema se forma subóptima... para, de alguna manera intentar optimizarlo. La idea que se propone en esta práctica es diferente de la ``del mono''. En el próximo apartado se os propone que se experimente con la función que actualiza la tabla $Q$ para llegar finalmente a la expresión de la ecuación de Bellman.



\subsection{El ciclo de \textit{train-test}}
El código que se proporciona tiene la siguiente idea:
\begin{itemize}
    \item Se ejecuta el script de train. Durante el entrenamiento se estima la matriz Q. Al final del entrenamiento esta matriz Q se guardará en el fichero \textt{qtable.npy}. 
    \begin{verbatim}
        qlearning.save_q_table(filename=’qtable.npy’)
    \end{verbatim}
     \item Realizado el entrenamiento, se deberá probar la tabla Q para ver su resultado. El script de test deberá cargar la matriz Q desde el fichero \textt{qtable.npy}. Así pues, se deberá modificar el script de test a esto:
     \begin{verbatim}
         from qlearning.qlearning import QLearning
         import gymnasium as gym

         def test_qlearning():
             environment = gym.make('CliffWalking-v1', render_mode="human")
             qlearning = QLearning(environment=environment)
             qlearning.load_q_table(filename='qtable.npy')
             print('Testing mode!')
             qlearning.test(total_episodes_test=1)

        if __name__ == "__main__":
             test_qlearning()
     \end{verbatim}
    



\subsection{Aprendizaje en el \textit{Cliff Walking environment}}
Abra el script \texttt{cliff\_q\_learning\_train.py}. Verá que es similar al que vimos anteriormente, con la diferencia de que en este caso guardamos la tabla Q al finalizar el método \emph{train}. La idea es: a) Entrenar y guardar Q. b) Testear: cargar Q y usarla para que el agente opere sobre el entorno. 

\begin{verbatim}
def train_qlearning():    
    environment = gym.make('CliffWalking-v1', render_mode='human')
    qlearning = QLearning(environment=environment)
    print('LET US LEARN NOW!')
    qlearning.train(total_episodes_train=300)    
    qlearning.save_q_table(filename='qtable.npy')

if __name__ == "__main__":
    train_qlearning()    
\end{verbatim}

Usa \texttt{render\_mode=human} si quieres ver la evolución del agente en el entorno. Normalmente, no haremos esto, pues ralentiza el aprendizaje:
\begin{verbatim}
    environment = gym.make('CliffWalking-v1', render_mode='human')
    environment = gym.make('CliffWalking-v1')
\end{verbatim}
    
Si entramos en el método train (clase Qlearning), veremos esto:
\begin{verbatim}
     def train(self, total_episodes_train):
        print("Training started!\n\n")

        
        return self.q_table  
\end{verbatim}

Como ayuda se proporciona el pseudo código que permite entrenar y calcular la tabla Q. La mayoría de los métodos ya se han visto en el método de test anterior.
\begin{verbatim}
    def train(self, total_episodes_train):
        for episode [in 1, N]:
            state = reset_environment()
            while True:
                action = take_random_action()
                next_state, reward, terminated = environment.apply_action(state, action)
                if terminated:                    
                    break
                update_q_table(state, next_state, action, reward)
                # Move to next state
                state = next_state
\end{verbatim}


\begin{exercise}
    Complete el método \texttt{train} de la clase \texttt{QLearning}. Nota: deje el método \texttt{update\_q\_table} vacío. Se completará en el apartado siguiente.   
\end{exercise}




\subsubsection{Actualización de Q (primera versión)}
En esencia, la tabla $Q(s_t,a_t)=r_t$ se puede entender de la siguiente manera:
\begin{itemize}
    \item Dado un estado $s_t$.
    \item Almacena un valor para cada una de las posibles acciones $a_t$ que puede ejecutar el agente en ese estado.
    \item Sabemos que el agente elegirá la acción como: $a_t=argmax(Q(s_t, a_t)).
\end{itemize}

Se propone una actualización de la tabla Q en la que se hace:
$$Q(s_t,a_t)=r_t$$

\begin{exercise}
    Complete la función \texttt{update\_q\_table} haciendo que:
    $$Q(s_t,a_t)=r_t$$
    Ejecute el algoritmo de \textit{train} varias veces y responda a las siguientes cuestiones:
    \begin{itemize}
        \item ¿El algoritmo consigue aprender?
        \item ¿El agente consigue completar la misión en algún momento?
    \end{itemize}
\end{exercise}


\subsubsection{¿Está aprendiendo el agente?}
Efectivamente, es necesario saber qué esta pasando en cada episodio durante el entrenamiento. En caso contrario, vamos ``a ciegas''. En el futuro nos enfrentaremos a problemas más complejos y es vital que sepamos si se está produciendo algún tipo de aprendizaje, pues, en caso contrario, deberíamos parar el algoritmo de aprendizaje y volver a la fase de diseño. 
\begin{exercise}
    \begin{itemize}
    \item Se propone que calcules la recompensa de cada episodio.
    \item Se propone que plotees (o lo imprimas por pantalla) la tabla Q al final del aprendizaje para ver qué está pasando. 
\end{itemize}
\end{exercise}



Según se ha podido comprobar, la actualización de la tabla Q anterior no es correcta. Así pues, la estrategia $Q(s_t, a_t)=r_t$ no produce buenos resultados. Debemos preguntarnos de nuevo: ¿qué representa $Q(s_t,a_t)$ en el problema? 

En el algoritmo Q-Learning, $Q(s_t,a_t)$ representa una estimación de todas las recompensas futuras que se obtendrían hasta el final del episodio al realizar la acción $a_t$ cuando el agente se encuentra en el estado $s_t$. Así pues,  $Q(s_t,a_t)=r_t$ es \textbf{inadecuada}, pues está almacenando solamente la recompensa actual y no la suma de recompensas hasta el final.

\subsubsection{Actualización de Q (segunda versión)}
Según se ha dicho, nos gustaría que $Q(s_t,a_t)$ guardara una estimación de las recompensas futuras que se obtendrían de estar en $s_t$ y tomar la acción $a_t$, es decir:
$Q(s_t,a_t)=r_t + r_{t+1} + r_{t+2} + r_{t+3}$. Para ello, vamos a considerar que tenemos un entorno donde el número máximo de pasos es $T$ (como ocurre en la mayoría de entornos de Gymnasium). Así pues:
\begin{equation}
Q(s_t,a_t) = \sum_{i=t}^{T} r_i= r_t +  r_{t+1}+ \ldots + r_T
\end{equation}
Suponga que en el instante siguiente, la matriz también debería almacenar la suma de recompensas hasta el final del episodio:
\begin{equation}
Q(s_{t+1},a_{t+1}) = \sum_{i=t+1}^{T} r_i = r_{t+1} +r_{t+2} + \ldots + r_T     
\end{equation}

Podemos hacer: 
\begin{equation}
    Q(s_t,a_t) - Q(s_{t+1},a_{t+1}) = r_t
\end{equation}
Con lo que tenemos otra forma de actualizar la tabla Q:
\begin{equation*}
Q(s_t,a_t) = r_t + Q(s_{t+1},a_{t+1})   
\end{equation*}

Durante la simulación, en un instante determinado, conocemos que el agente está en $s_t$ y realiza la acción $a_t$. Como consecuencia, en ese mismo instante, el agente recibirá una recompensa $r_t$ (por ejemplo, si el agente se cae por el acantilado, $r_t=-100$. Si no se cae: $r_t=-1$). El problema ahora es que debemos sumar el término $Q(s_{t+1},a_{t+1})$. De este término, conocemos $s_{t+1}$ y tenemos 4 valores de Q (uno para cada acción posible). Con esto, ¿qué valor elegimos de $Q(s_{t+1},a_{t+1})$?. Esta parte se soluciona eligiendo el máximo de estos valores (). Así, se propone una \textbf{forma mejorada de la ecuación de actualización de la tabla Q}:
\begin{equation*}
Q(s_t,a_t) = r_t + max_a(Q(s_{t+1},a))
\end{equation*}

Nótese que esta ecuación se debe ejecutar de forma repetida para poder llegar a una estimación final, pues inicialmente se desconoce la tabla Q. 

\subsubsection{El factor de descuento $\gamma$}
En algunos entornos, la suma de recompensas siguiente no permite estimar correctamente la matriz. En la literatura se propone modificar la suma de recompensas:
\begin{equation}
Q(s_t,a_t) = \sum_t^{T} r_i= r_t +  r_{t+1}+ \ldots + r_T
\end{equation}
Por esta suma de recompensas:
\begin{equation}
Q(s_t,a_t) = r_t + \gamma r_{t+1}+ \gamma^2 r_{t+2} + \gamma^3r_{t+3} +  \ldots + \gamma^T r_T = \sum_{i=t}^{T} \gamma^{i-1}r_i
\end{equation}
Con este factor $0\leq\gamma\leq1$ se consigue que el agente pueda ponderar más las recompensas cercanas al instante $t$ (los pasos siguientes). Igual que hicimos antes:
\begin{eqnarray*}
Q(s_t,a_t) &=& r_t + \gamma r_{t+1}+ \gamma^2 r_{t+2} + \gamma^3r_{t+3} + \gamma^4r_{t+4} +  \ldots \\
Q(s_{t+1},a_{t+1}) &=& r_{t+1} + \gamma r_{t+2}+ \gamma^2 r_{t+3} + \gamma^3r_{t+4} + \gamma^4r_{t+5} +  \ldots \\
\end{eqnarray*}

Aquí, vemos que se puede escribir:
\begin{eqnarray*}
Q(s_t,a_t) &=& r_t + \gamma 
Q(s_{t+1},a_{t+1}) 
\end{eqnarray*}
E, igual que antes, nos deja la actualización siguiente de la tabla Q:
\begin{eqnarray*}
Q(s_t,a_t) &=& r_t + \gamma 
\max_a(Q(s_{t+1},a))
\end{eqnarray*}

Obviamente, si ponemos $\gamma=1$, tenemos el caso en el que el agente suma todas las recompensas sin ponderar. Si ponemos, por ejemplo, $\gamma=0.5$, el agente considera solamente las recompensas cerca del instante actual y es bastante ``miope'' en relación con las recompensas futuras.

\subsubsection{El factor de aprendizaje $\alpha$}
Siempre que planteamos un proceso de aprendizaje, deberíamos plantear una forma de actualización del valor actual de un parámetro. Nótese que en la última ecuación se hacía:
\begin{equation*}
Q(s_t,a_t) = r_t + \gamma 
max_a(Q(s_{t+1},a_{t+1}))
\end{equation*}
Con lo cual, el valor de $Q(s_t,a_t)$ se sustituye directamente con el nuevo valor $Q(s_t,a_t) = r_t + \gamma 
\max_a(Q(s_{t+1},a_{t+1}))$. En el ámbito de la inteligencia artificial, generalmente se propone utilizar un factor de aprendizaje $\alpha$ que permite añadir más o menos información a la estimación actual en cada paso del algoritmo. Imaginemos que en la iteración $i$ tenemos una estimación $Q(s_t,a_t)_i$ de una celda de la tabla. La nueva estimación en la iteración siguiente es $Q(s_t,a_t)_{i+1}$. La nueva información de la que se dispone en el instante $i$ es $\hat{Q}(s_t,a_t) = r_t + \gamma 
max_a(Q(s_{t+1},a_{t+1}))$:
\begin{equation}\label{eq:learning_rate}
Q(s_t,a_t)_{i+1} = (1-\alpha) Q(s_t,a_t)_{i} + \alpha \hat{Q}(s_t,a_t)
\end{equation}
Con este parámetro $\alpha$ podemos regular la forma en la que se añade nuevo aprendizaje a la tabla Q. Así pues:
\begin{itemize}
    \item Si $\alpha=1$ el aprendizaje es máximo, y el algoritmo olvida la información anterior en cada instante de simulación.
    \item Si $\alpha=0$ el aprendizaje es mínimo y no se añade nueva información a la tabla Q.
\end{itemize}

\section{La ecuación de Bellman}
Con todo lo anterior, escribimos la ecuación de actualización completa sustituyendo en la ecuación (\ref{eq:learning_rate}) el valor de actualización propuesto $\hat{Q}(s_t,a_t) = r_t + \gamma 
\max_a(Q(s_{t+1},a_{t+1}))$. Quedando:
\begin{equation}
Q(s_t,a_t)_{i+1} = Q(s_t,a_t)_{i} + \alpha [r_t + \gamma 
\max_a(Q(s_{t+1},a))- Q(s_t,a_t)_{i}]    
\end{equation}

Así pues, esta última ecuación se conoce como la Ecuación de Bellman. El matemático Richard Ernest Bellman la propuso en los años 50. La ecuación nos deja dos parámetros muy interesantes ($\alpha$ y $\gamma$) que podemos variar para resolver diferentes problemas de Q-learning.

Así pues, en resumen, tenemos los siguientes parámetros

\section{Resumen}
El código siguiente permite resolver un problema de aprendizaje e incluye todos los conceptos que hemos visto
\begin{verbatim}
    
    def train(self, total_episodes):
        print("Training started!\n\n")
        # Training loop
        for episode in range(total_episodes):
            print("Episode: ", episode)
            state, info = self.env.reset()
            # sum of rewards for each episode
            srt = 0
            while True:
                # Epsilon-greedy action selection
                if random.uniform(0, 1) < self.epsilon:
                    action = self.env.action_space.sample()  # Explore
                else:
                    action = np.argmax(self.q_table[state])  # Exploit
                # Take action, observe new state and reward
                next_state, reward, terminated, truncated, info = self.env.step(action)
                srt += reward
                if terminated or truncated:
                    self.results.save_data(sum_rewards_per_episodei=srt)
                    # self.results.q_table = self.q_table
                    break
                # Actualiza la tabla
                self.update_q_table(state, action, next_state, reward)
                # Move to next state
                state = next_state
            # Reduce epsilon (less exploration, more exploitation as time goes on)
            # self.epsilon = self.min_epsilon + (self.max_epsilon - self.min_epsilon) * np.exp(-self.decay_rate * episode)
            self.epsilon *= self.decay_rate
            self.epsilon = max(self.epsilon, self.min_epsilon)
        print("Training finished! Your Q-table is optimized.")
        self.env.close()
        return self.q_table


    def   update_q_table(self, state, action, next_state, reward):
        #  Update Q-table
        # self.q_table[state, action] = reward
        #self.q_table[state, action] = reward + np.max(self.q_table[next_state])
        # Step 3: Update Q-table using the Bellman Equation
         self.q_table[state, action] = (self.q_table[state, action] +
                                        self.learning_rate * (
                                                    reward + self.discount_rate * np.max(self.q_table[next_state]) -
                                                    self.q_table[state, action]))

    def test(self, total_episodes):
        print('Test started')
        # test loop
        for episode in range(total_episodes):
            print("Episode: ", episode)
            state, info = self.env.reset()
            srt = 0
            while True:
                # Step 1: Greedy action selection
                action = np.argmax(self.q_table[state])  # Exploit
                # Step 2: Take action, observe new state and reward
                next_state, reward, terminated, truncated, info = self.env.step(action)
                srt += reward
                time.sleep(.1)
                if terminated or truncated:
                    self.results.save_data(sum_rewards_per_episodei=srt)
                    break
                # Move to the next state
                state = next_state
        print("Test finished!")
        self.env.close()
\end{verbatim}
%En resumen, hemos visto cómo actualizar la tabla de acuerdo con la ecuación de Bellman. 
%\begin{equation}
%Q(s_t,a_t)_{i+1} = Q(s_t,a_t)_{i} + \alpha [r_t + \gamma 
%\max_a(Q(s_{t+1},a))- Q(s_t,a_t)_{i}]    
%\end{equation}






\section{Otro entorno más desafiante: Taxi v4}

En el caso del taxi, las recompensas se definen así:
Rewards
-1 per step unless another reward is triggered.
+20 delivering passenger.
-10 executing “pickup” and “drop-off” actions illegally.
Así pues, el hecho de realizar un paso (step) conlleva una recompensa negativa (), pues buscamos que el taxi realice el transporte en el mínimo tiempo posible. Cuando se deja al pasajero se obtiene una recompensa positiva (+20). Nótese que se obtiene una recompensa negativa por acciones ilegales de recogida y reparto de la persona: pues podría ocurrir que el agente decidiera dejar al pasajero lejos de su destino o bien intentar recogerlo en una casilla donde no estuviera la persona).

En el caso del taxi, los estados iniciales serán diferentes. Esto incluye variaciones de la posición inicial del taxi, de la posición inicial del pasajero y del destino. Se excluyen de estos estados iniciales (la librería gymnasium los elimina) los casos ilegales en los que el pasajero ya se encuentra dentro del vehículo, el vehículo y el pasajero se encuentran ya en el destino… etc.


Taxi v4
Explicar estados. Codificación.
Acciones.
Recompensas.
Cliff
Frozen lake.
Más sencillo que el caso del Taxi, pero en este caso no se trata de un proceso determinista (aunque sí un proceso de Markov) p(s’|s, a)= no
Estados
Acciones
Recompensas

Descripción de cada uno de los entornos

Taxi v4 descripción
Starting State
The initial state is sampled uniformly from the possible states where the passenger is neither at their destination nor inside the taxi. There are 300 possible initial states: 25 taxi positions, 4 passenger locations (excluding inside the taxi) and 3 destinations (excluding the passenger’s current location).
Rewards
-1 per step unless another reward is triggered.
+20 delivering passenger.
-10 executing “pickup” and “drop-off” actions illegally.
An action that results a noop, like moving into a wall, will incur the time step penalty. Noops can be avoided by sampling the action_mask returned in info.
Episode End
The episode ends if the following happens:
Termination:
The taxi drops off the passenger.
Truncation (when using the time_limit wrapper):
The length of the episode is 200.
Q learning, descripción funcionamiento

Vamos a resolver los tres entornos, que se pueden considerar ordenados en dificultad


\section{Instalación}
Los equipos de prácticas tienen todo lo necesario para realizar la práctica. En este apartado se dan algunas instrucciones para que puedas instalar todo lo necesario en otro portátil/PC y poder trabajar de forma autónoma. Las instrucciones se ofrecen para un sistema Linux/Ubuntu. 
Necesitamos instalar el programa git
\begin{verbatim}
$ sudo apt install git
\end{verbatim}

Comenzamos por descargarnos el repositorio de prácticas. Es un repositorio que se encuentra en Github:
\begin{verbatim}
$ cd /home/usuario/Desktop
$ git clone https://github.com/4rtur1t0/Intro2AI
\end{verbatim}
Veremos que tenemos un directorio llamado Intro2AI en el escritorio. Seguimos instalando algunos paquetes necesarios para la creación de entornos virtuales de python \texttt{virtualenv}:
\begin{verbatim}
sudo apt install python3-pip python3-venv python3-dev -y 
\end{verbatim}

Seguidamente, creamos un entorno virtual:
\begin{verbatim}
$ cd /home/usuario
$ mkdir Applications
$ cd Applications
$ virtualenv venvAI
$ source /home/usuario/Applications/venvAI/bin/activate
\end{verbatim}
Con esto hemos creado un entorno virtual de python en el directorio \texttt{/home/usuario/Applications/venvAI} y lo hemos activado. Instalamos ahora algunos paquetes para este entorno de python:
\begin{verbatim}
$ cd /home/usuario/Applications/venvAI/bin
$ ./pip install -r /home/arturo/Desktop/Intro2AI/requirements.txt 
$ ./pip install stable-baselines3 swig torch
$ ./pip install tensorboard
$ ./pip install "gymnasium[atari,accept-rom-license]" stable-baselines3 shimmy
\end{verbatim}

Cuando usemos un editor de python, como, pyCharm o Visual Studio, necesitaremos asociar el entorno virtual que hemos creado con el proyecto.

import gymnasium as gym
import numpy as np
import random

# Initialize environment
env = gym.make("Taxi-v3")

# Initialize Q-table with zeros (500 states x 6 actions)
state_size = env.observation_space.n
action_size = env.action_space.n
q_table = np.zeros((state_size, action_size))

# Hyperparameters
total_episodes = 2000
learning_rate = 0.1      # alpha
discount_rate = 0.6      # gamma
epsilon = 1.0            # Exploration rate
max_epsilon = 1.0
min_epsilon = 0.01
decay_rate = 0.01        # Exponential decay rate for exploration

# Training loop
for episode in range(total_episodes):
    state, info = env.reset()
    done = False
    
    while not done:
        # Step 1: Epsilon-greedy action selection
        if random.uniform(0, 1) < epsilon:
            action = env.action_space.sample() # Explore
        else:
            action = np.argmax(q_table[state]) # Exploit
            
        # Step 2: Take action, observe new state and reward
        next_state, reward, terminated, truncated, info = env.step(action)
        done = terminated or truncated
        
        # Step 3: Update Q-table using the Bellman Equation
        q_table[state, action] = q_table[state, action] + learning_rate * (
            reward + discount_rate * np.max(q_table[next_state]) - q_table[state, action]
        )
        
        # Move to next state
        state = next_state
        
    # Reduce epsilon (less exploration, more exploitation as time goes on)
    epsilon = min_epsilon + (max_epsilon - min_epsilon) * np.exp(-decay_rate * episode)

print("Training finished! Your Q-table is optimized.")
env.close()


Cuestiones que se pueden preguntar:
Variaciones sobre estos parámetros. ¿En qué afectan? → a la memoria de prácticas.
total_episodes = 2000
learning_rate = 0.1      # alpha
discount_rate = 0.6      # gamma
epsilon = 1.0            # Exploration rate
max_epsilon = 1.0
min_epsilon = 0.01
decay_rate = 0.01    
¿qué pasa cuando variamos el decay rate? ¿qué se observa? ¿qué pasa si se hace una función lineal o exponencial?

La solución es óptima (la trayectoria es óptima).

Suponga que el pasajero está a bordo. Para todos los estados (todas las casillas).
Haga un mapa que muestre el valor de la acción SOUTH sobre todas las casillas
Haga un mapa que muestre el valor de la acción WEST sobre todas las casillas
Haga un mapa que muestre el valor de la acción SOUTH-WEST sobre todas las casillas.


Podemos seguir aprendiendo mientras se opera en modo funcionamiento. En nuestro problema del taxi se opera en una ciudad. Imagine un taxista experimentado que trabaja en Madrid. Un buen día (o malo) le cortan la castellana (una de las casillas ). ¿Puede el taxista aprender esta situación?



¿FUNCIONA?
Pues en efecto, se os proporciona poco código y estamos luchando para que lo escribáis vosotros mismos.
Pues le puedes preguntar a la IA si funciona tu código, aquí tienes una probabilidad de que te dé una versión probada del algoritmo de los muchos repositorios que hay por internet de gente que ha encontrado soluciones.
El objetivo es que tú aprendas a encontrar estas soluciones y a encontrar


EJERCICIO
Nos gustaría poder analizar el reward de cada episodio, para las acciones que se tomen de acuerdo con el conocimiento de la tabla… y no las que se tomen aleatoriamente. →

Podríamos forzar algún episodio durante el entrenamiento con epsilon nulo, para evaluar el desempeño.

EJERCICIO
Evalúe el entrenamiento en base a la recompensa media.







\bibliographystyle{alpha}
\bibliography{sample}

\end{document}
